\slajdid{02-s005}
\begin{frame}[label=02-s005]{Maksimalna specificirana kašnjenja}
\begin{columns}[T,onlytextwidth]
\begin{column}{.39\textwidth}\centering
% element: 02-s005:e01
\begin{tikzpicture}[DEoznaka,x=8.6mm,y=8mm]
\draw[DEstrelica] (-.7,2.25)--(3.8,2.25) node[below] {\(t\)};
\draw[DEstrelica] (0,2.10)--(0,4.05) node[left] {\(x(t)\)};
\draw[DEstrelica] (-.7,-.1)--(3.8,-.1) node[below] {\(t\)};
\draw[DEstrelica] (0,-.25)--(0,1.7) node[left] {\(y(t)\)};
\node[above left,inner sep=1pt] at (0,3.35) {=1};\node[above left,inner sep=1pt] at (0,2.45) {=0};
\node[above left,inner sep=1pt] at (0,1) {=1};\node[above left,inner sep=1pt] at (0,.1) {=0};
\draw[DEvod] (-.5,2.45)--(.6,2.45)--(.6,3.35)--(1.6,3.35)--(1.6,2.45)--(3.5,2.45);
\draw[DEvod] (-.5,1)--(1.15,1)--(1.15,.1)--(2.5,.1)--(2.5,1)--(3.5,1);
\fill[DEcrvena!15] (.6,.1) rectangle (1.15,1);
\fill[DEcrvena!15] (1.6,.1) rectangle (2.5,1);
\foreach \a/\b in {.6/1.15,1.6/2.5}{
 \begin{scope}
 \clip (\a,.1) rectangle (\b,1);
 \foreach \x in {0,.12,...,3.5}{\draw[DEcrvena,thin] (\x,.1)--++(.9,.9);}
 \end{scope}
 \draw (\a,.1) rectangle (\b,1);
}
\draw[dashed] (0,3.35)--(.6,3.35) (0,.1)--(.6,.1);
\draw[dashed] (.6,2.45)--(.6,-.3) (1.6,2.45)--(1.6,-.3);
\draw[dashed] (1.15,.1)--(1.15,-.3) (2.5,.1)--(2.5,-.3);
\draw[DEcrvena,<->] (.6,-.25)--(1.15,-.25) node[midway,below] {\(t_{pHL}\)};
\draw[DEcrvena,<->] (1.6,-.25)--(2.5,-.25) node[midway,below] {\(t_{pLH}\)};


\draw[DEcrvena,-{Latex[length=1.6mm]}] (.65,-1.3)--(1.15,-.8);
\draw[DEcrvena,-{Latex[length=1.6mm]}] (3.1,-1.3)--(2.5,-.8);
\node[DEcrvena,below,inner sep=1pt] at (.65,-1.3) {Sigurno 0};
\node[DEcrvena,below,inner sep=1pt] at (3.1,-1.3) {Sigurno 1};
\end{tikzpicture}

\par\vspace{2mm}\Oznaka{Šrafirano: izlaz može biti 0 ili 1.}
\end{column}
\begin{column}{.58\textwidth}
\Rezultat{\alert{Maksimalna} vremena \(t_{pHL}\) i \(t_{pLH}\) određuju kada smo \alert{sigurni} u novo stanje izlaza.}
\par\vspace{2mm}
Promena može nastupiti i ranije; po isteku odgovarajućeg maksimalnog vremena novo stanje je sigurno uspostavljeno.
\par\vspace{2mm}
% element: 02-s005:e02
Stvarno kašnjenje zavisi od napajanja, temperature, starosti i konkretnog primerka kola: \alert{promenljivo je i za isto kolo}.
\par\vspace{2mm}
% element: 02-s005:e03
Proizvođač specificira granice za \alert{tip kola}. One važe za \alert{sve primerke} tog tipa u specificiranim uslovima; stvarna kašnjenja primeraka mogu se razlikovati.
\end{column}
\end{columns}
\end{frame}
